\documentclass[10pt,a4paper]{amsart}
\usepackage[margin=19mm]{geometry}
\usepackage{amsmath,amssymb,mathtools}
\usepackage[scr=rsfs,cal=euler]{mathalfa}
\usepackage{xfrac}
\usepackage[unicode,hidelinks,bookmarks=false]{hyperref}
\newcommand{\ie}{i.e.\ }
\newcommand{\cf}{cf.\ }
\newcommand{\resp}{respectively\ }
\newcommand{\Subsection}{\subsection}
\providecommand{\nobreakdash}{\nobreak\hbox{-}}
\setlength{\emergencystretch}{2em}
\multlinegap0pt
\usepackage{amssymb}
\usepackage{bm}
\usepackage{mathtools}
\usepackage{enumerate}
\newtheorem{lemma}{Lemma}
\newtheorem{proposition}{Proposition}
\DeclareMathOperator{\Sym}{Sym}
\title{Counting parabolic principal $G$-bundles with nilpotent sections over $\mathbb{P}^1$}
\author{Rahul Singh}
\date{Source: arXiv:2111.11583v2 (CC BY 4.0)}
\begin{document}
\maketitle

\noindent\emph{Source and licence.} Counting parabolic principal $G$-bundles with nilpotent sections over $\mathbb{P}^1$. Version arXiv:2111.11583v2 (CC BY 4.0), distributed by arXiv under the Creative Commons Attribution 4.0 International licence, http://creativecommons.org/licenses/by/4.0/, as recorded in the arXiv licence metadata for this exact version and shown on the abstract page https://arxiv.org/abs/2111.11583v2. Full licence notice: \texttt{licenses/arxiv-2111.11583-CC-BY-4.0.txt}.\par\smallskip

\noindent\textit{Editorial scope.} Counted unit: the article's proposition coproduct (source0.tex lines 1628--1633). The article's complete original proof of it is delivered with it (source0.tex lines 1634--1681, one \texttt{proof} environment of 48 lines). No mathematics is written, repaired or re-derived: the statement and the proof are the article's own lines, in the article's own notation, and no retained line is substituted or re-flowed.  Retained uncounted as the source-local context the proof needs: lines 698--712.  Nothing was omitted from any retained span, so the package carries no omission record.  No retained line contains a citation, so the wrapper carries no bibliography and no bibliography entry.  No retained line contains a figure environment, an image-inclusion command or drawing code, so no image is delivered and the package declares no asset dependency.  Editorial formatting only, in addition: an AMS \texttt{amsart} preamble that restates the article's own declarations and macros used by the retained lines, namely the environments \texttt{lemma}, \texttt{proposition} on separate counters, together with the standard packages those lines need.  The retained environments print in this document as Lemma 1, Proposition 1 in order of appearance, whereas the article prints the same items with the numbering of its own full document.  The complete native source of the article as distributed in the arXiv e-print for this exact version is bundled unchanged as \texttt{sources/118768/source0.tex}.
\begin{lemma}\label{associateinvariantdelta}
%Keep notations as above. Then
We have
\begin{equation}\label{lhsrhs}
\Delta_{H}(\delta_{\mathcal{O}})=\sum_{(\mathcal{O}_{1},\mathcal{O}_{2})\in\big(\mathcal{P}(\Pi_{H})/\sim\big)\times\big(\mathcal{P}(\Pi_{H})/\sim\big)}n_{\mathcal{O}}^{\mathcal{O}_{1},\mathcal{O}_{2}}\delta_{\mathcal{O}_{1}}\otimes \delta_{\mathcal{O}_{2}},
\end{equation}
where
\[
n_{\mathcal{O}}^{\mathcal{O}_{1},\mathcal{O}_{2}}
=
\big|\{w\in W_{J_{\mathcal{O}_{1}}}\backslash W_{H}/W_{J_{\mathcal{O}_{2}}}:\Phi_{J_{\mathcal{O}_{1}}}\cap w\cdot \Phi_{J_{\mathcal{O}_{2}}}=
w'\cdot\Phi_{J_{\mathcal{O}}}\text{ for some }w'\in W_H\}\big|.
\]
In particular, $\Delta_H$ preserves the subspace of associate invariant functions.
\end{lemma}

\begin{proposition}\label{coproduct agrees with coproduct}
We have
$
\Delta'_{GL_{n}}=\Delta^n.
$
\end{proposition}

\begin{proof}
Since $m_{\nu}$, $\nu\in\Pi_{n}$ form a basis of $\Sym_{\mathbb{Z}}^{n}[X]$, it is enough to check that $\Delta'_{GL_{n}}$ agrees with $\Delta^{n}$ on this basis.
We re-write the conclusion of Lemma \ref{associateinvariantdelta} for $GL_n$. Let $\nu$ be a partition of $n$. It gives an equivalence relation $\sim_\nu$ on $\{1,\ldots,n\}$, where $i\sim_\nu j$ if and only if there exists a $t$ such that $\nu_1+\ldots+\nu_t\leq i,j<\nu_1+\ldots+\nu_{t+1}$. Note that $S_n$ acts on $\{1,\ldots,n\}$ and thus on the equivalence relations. For an equivalence relation $\sim$ on $\{1,\ldots,n\}$, we will write $Part(\sim)\in\Pi_n$ for the corresponding partition of $n$, that is, the ordered sequence of sizes of equivalence classes. 
By Lemma \ref{associateinvariantdelta}, we have
\[
    \Delta'_{GL_{n}}(m_\nu)=\sum_{\lambda,\mu\in\Pi_{n}}n_{\nu}^{\lambda,\mu}m_\lambda\otimes m_\mu,
\]
where 
\[
n_{\nu}^{\lambda,\mu}=\big|\{w\in S_{\lambda}\backslash S_{n}/S_{\mu}:Part(\sim_{\lambda}\cap w(\sim_{\mu}))=\nu\}\big|.
\]
%If we identify roots for $GL_n$ with pairs of integers, then $\Phi(J_\nu)$ is equal to $\sim_\nu$ without the diagonal.
Thus we have
\begin{equation}\label{eq:GLn}
  n_\nu^{\lambda,\mu}
  =
  \sum_{\substack{w\in S_n\colon
  \\
  Part(\sim_\lambda\cap w(\sim_\mu))=\nu}}\frac{|w^{-1}S_{\lambda}w\cap S_{\mu}|}{|S_\mu||S_\lambda|}.
\end{equation}
Now consider the coproduct $\Delta^n$.
We have
\[
    \Delta^{n}(m_\nu)=\sum_{[(i_1,j_1),\ldots,(i_n,j_n)]}(X_{i_1}Y_{j_1})\ldots(X_{i_n}Y_{j_n})
\]
where the sum is over all multisets $[(i_t,j_t)]$ such that the multiplicities of elements are given by $\nu$.

The group $S_n$ is acting naturally on length $n$ sequences. Let $(\mu)$ be the standard sequence
\[
    \underbrace{1,\ldots,1}_{\mu_1\text{ times}},\underbrace{2,\ldots,2}_{\mu_2\text{ times}},\ldots.
\]
In $\Delta^{n}(m_{\nu})$, $X^\lambda Y^\mu$ occurs as:
\[
    \sum_{j_1,\ldots,j_n}\frac{1}{|\text{orbit of }S_\lambda\text{ on }j_{1},\ldots,j_{n}|}(X_1Y_{j_1}\ldots X_1Y_{j_{\lambda_1}})(X_2Y_{j_{\lambda_1+1}}\ldots X_2Y_{j_{\lambda_2}})\ldots,
\]
where the summation is over all sequences $j_1,\ldots,j_n$ such that $[j_{1},\ldots,j_{n}]=[(\mu)]$ and $Part(\sim_\lambda\cap\sim_j)=\nu$, where $\sim_j$ denotes the equivalence relation $t\sim_j s$ iff $j_t=j_s$.
Let $\widetilde{n}_{\nu}^{\lambda,\mu}$ denote the coefficient of $m_{\lambda}\otimes m_{\mu}$ in $\Delta^{n}(m_{\nu})$.
The condition $[j_1,\ldots,j_n]=[(\mu)]$ is equivalent to the existence of $w\in S_n$ such that $w\cdot(\mu)=(j_1,\ldots,j_n)$, in which case $w(\sim_\mu)=\sim_j$. Since there are exactly $|S_\mu|$ such $w$, we get
\[
\widetilde{n}_\nu^{\lambda,\mu}
  =
  \sum_{\substack{w\in S_n\colon
  \\
  Part(\sim_\lambda\cap w(\sim_\mu))=\nu}}\frac{|w^{-1}S_{\lambda}w\cap S_{\mu}|}{|S_\mu||S_\lambda|},
\]
which agrees with~\eqref{eq:GLn}.
This finshes the proof of Proposition \ref{coproduct agrees with coproduct}.
\end{proof}

\end{document}
